% TOP_PROBLEM: {"id": 2301026, "title": "Research Problems in Function Theory — Problem 1.26", "category_id": 8, "category": {"id": 8, "name": "analysis", "display_name": "Analysis", "description": "Limits, continuity, calculus, and function theory.", "slug": "analysis", "order_index": 8, "created_at": "2026-07-31T15:26:25.671Z"}, "status": "partially_solved", "problem_number": "AMR-022-1026", "set_id": 13}
% FIRST_POSTED: 2026-10-02
% ENTRY_KIND: statement_only
% UnsolvedMath ID 2301026; source catalogue code AMR-022-1026
% !TeX program = lualatex
% Definition and sources only; this file is not an LLM solution attempt.
\documentclass[11pt,a4paper]{article}
\usepackage[margin=25mm]{geometry}
\usepackage{fontspec}
\setmainfont{lmroman10-regular.otf}[BoldFont=lmroman10-bold.otf,
 ItalicFont=lmroman10-italic.otf,BoldItalicFont=lmroman10-bolditalic.otf]
\usepackage{amsmath,amssymb,mathtools}
\usepackage{unicode-math}
\setmathfont{latinmodern-math.otf}
\AtBeginDocument{\let\setminus\smallsetminus}
\usepackage{xurl,hyperref}
\hypersetup{colorlinks=true,urlcolor=blue}
\setcounter{secnumdepth}{0}
\setlength{\parindent}{0pt}
\setlength{\parskip}{5pt}
\setlength{\emergencystretch}{3em}
\begin{document}
\section*{Research Problems in Function Theory — Problem 1.26}\label{2301026}
{\small UnsolvedMath ID 2301026 · AMR-022-1026 · Analysis · AMR Open Problem Lists}
\subsection{Definitions and mathematical statement}
For a nonconstant meromorphic function $f$, let $n(t,a)$ count solutions of $f(z)=a$ in $|z|\le t$, with multiplicity; $a=\infty$ means poles. Put
\[
N(r,a)=\int_0^r\frac{n(t,a)-n(0,a)}t\,dt+n(0,a)\log r,
\quad T(r,f)=N(r,\infty)+\frac1{2\pi}\int_0^{2\pi}\log^+|f(re^{i\theta})|\,d\theta,
\]
where $\log^+x=\max(0,\log x)$. At a boundary radius $R\in\{1,\infty\}$ with $T(r,f)\to\infty$, define the deficiency
$\delta(a,f)=1-\limsup_{r\to R}N(r,a)/T(r,f)$. A deficient value has positive deficiency.
Suppose $f$ is meromorphic in the plane of finite order $\rho>1$, defined by $\limsup\log T(r,f)/\log r$. Set $q=\lfloor2\rho\rfloor$ and $s=\frac\pi2(2\rho-q)$. The Drasin--Weitsman conjecture asks whether
\[
\sum_{a\in\widehat{\mathbb C}}\delta(a,f)
\le \max\left\{2-\frac{2\sin s}{q+2\sin s},\quad
2-\frac{2\cos s}{q+1}\right\}.
\]
The sum means the supremum of sums over finite collections of distinct values, so it is meaningful without first enumerating them. The proposed inequality refines the usual bound two according to the order of growth. The source states that the candidate bound is sharp and that the range $0\le\rho\le1$ is settled. The task here is the universal inequality in the remaining finite-order range, including the half-integral endpoints under the displayed convention for $q$.
\subsection{Short English statement}
Prove the proposed sharp bound, depending on growth order, for the total deficiency of a meromorphic function.
\subsection{Sources}
\begin{itemize}
\item[S1] ULAM.AI and the UnsolvedMath Contributors, supplied problems.json, record 2301026 (AMR-022-1026), AMR Open Problem Lists. Catalogue identity and source transcription; CC BY 4.0.
\url{https://huggingface.co/datasets/ulamai/UnsolvedMath}.
\item[S2] Hayman - Research Problems in Function Theory (2018), Problem 1.26. Original source indexed by AMR-022-1026.
\url{https://arxiv.org/abs/1809.07200}.
\end{itemize}
\end{document}
